\section{Experimental Results}
\label{section:exp}
\subsection{Evaluation of Baseline Methods}
\label{section:exp:baseline}

At the current proposal stage, we are actively constructing our custom synthetic dataset (Staff-to-Jianpu image-sequence pairs) and configuring the baseline models. We plan to conduct our preliminary experiments by evaluating the classical Image-to-Sequence baseline (CRNN) and the modern vision-language baseline (TrOCR) on our generated validation set.

To rigorously assess the baseline performance, we will report the Symbol Error Rate (SER) and the Sequence-level Exact Match Rate. The expected format for our preliminary baseline evaluation is presented in Table \ref{table:baseline_results}. 

\begin{table}[htbp]
\centering
\caption{Preliminary evaluation of baseline methods on the synthetic validation set.}
\label{table:baseline_results}
\begin{tabular}{lcc}
\toprule
\textbf{Method} & \textbf{Symbol Error Rate (SER) $\downarrow$} & \textbf{Exact Match Rate $\uparrow$} \\
\midrule
Baseline 1: Oemer (with heuristic converter) & TBD & TBD \\
Baseline 2: CRNN (ResNet + RNN + CTC) & TBD & TBD \\
Baseline 3: TrOCR (ViT + Transformer) & TBD & TBD \\
\bottomrule
\end{tabular}
\end{table}

These preliminary results will serve as a foundation for designing and tuning our proposed end-to-end architecture. We anticipate that the Transformer-based baseline (TrOCR) will outperform the CTC-based baseline (CRNN) due to its superior sequence modeling capability, which is crucial for capturing the underlying musical syntax and alignment rules without explicit programming.